\section{Why a finite angular ansatz cannot work}
\label{app:app-obstructions}

This appendix is independent of the existence proof. It explains one
feature of the construction: although the seed $\Phi_*$ has the single
angular frequency two, the inverse of the linearized operator had to
be constructed for \emph{all} even frequencies.
Proposition~\ref{prop:obs-monotone-tail} ruled out a monotone flat end
for the linearized equation. Another tempting simplification would be
to look for the nonlinear height among functions with finitely many
angular modes. The following proposition shows that this is
impossible as soon as the height is required to be radial on a
collar.

\begin{proposition}\label{prop:obs-finite-square}
Let $E>0$ be a $C^1$ function on an interval $I\subset(0,\infty)$,
and let $z$ be a real $C^2$ solution of the Darboux equation
\eqref{eq:intro-darboux} for the metric $E\,dr^2+r^2\,d\theta^2$ on
$I\times\T$. Suppose that
\[
 z^2=\sum_{|j|\le J}b_j(r)\,e^{ij\theta},\qquad
 b_{-j}=\overline{b_j},\qquad b_j\in C^2(I),
\]
and that $z$ is independent of $\theta$ for $r$ in some nonempty open
subinterval of $I$. Then $z$ is independent of $\theta$ on all of
$I$.
\end{proposition}

The hypothesis covers every height with finitely many angular modes,
and every smooth square root of a trigonometric polynomial in
$\theta$.

The difficulty in the proof is the continuation through a zero of a
Fourier coefficient. The highest coefficient satisfies an ordinary
differential equation which is singular where the coefficient
vanishes, and the coefficient is complex, so neither the standard
uniqueness theorem nor an argument for real functions applies
directly.

\begin{proof}
\step{1: a polynomial identity for $v=z^2$}
Let $K=E'/(2rE^2)$ be the curvature of the metric. From
$\nabla^2v=2z\,\nabla^2z+2\,dz\otimes dz$ and the expansion of a
$2\times2$ determinant under a perturbation of rank one,
\begin{equation}\label{eq:obstruction-polynomial}
 2v\,\det\nolimits_g\nabla^2v
 -\operatorname{cof}_g(\nabla^2v)[dv,dv]
       =8Kv^2-2Kv\,|dv|_g^2 .
\end{equation}
Indeed, on the left the terms containing
$z^3\operatorname{cof}_g(\nabla^2z)[dz,dz]$ cancel, leaving
$8z^4\det_g\nabla^2z$, and the Darboux equation gives the right side.
No division by $z$ was used, so the identity holds also where $z=0$.

\step{2: the equation for the highest coefficient}
Let $b=b_J$ be the coefficient of the highest frequency $J>0$, and put
$a=EJ^2/r$ and $d=E'J^2/(2r)$. In \eqref{eq:obstruction-polynomial}
the frequency $3J$ arises only from products of three copies of
$b\,e^{iJ\theta}$. Write the entries of the Hessian of $v$ as
\[
 A=v_{rr}-\frac{E'}{2E}\,v_r,\qquad
 B=v_{r\theta}-\frac{v_\theta}r,\qquad
 C=v_{\theta\theta}+\frac rE\,v_r ,
\]
so that, after multiplication by $Er^2$, the left side of
\eqref{eq:obstruction-polynomial} is
$2v(AC-B^2)-(Cv_r^2-2Bv_rv_\theta+Av_\theta^2)$. Replacing
$v,v_r,v_{rr},v_\theta,v_{r\theta},v_{\theta\theta}$ by
$b,b',b'',iJb,iJb',-J^2b$ and including the cubic term of the right
side, the coefficient of $e^{3iJ\theta}$ gives
\begin{equation}\label{eq:obstruction-highest-mode}
\begin{split}
0={}&(2bb'-ab^2)\,b''-(b')^3+ab\,(b')^2
 -\frac{2a}{r}\,b^2b'+\frac{2a}{r^2}\,b^3\\
 &\quad+d\,b^2b'-\frac{2d}{r}\,b^3 .
\end{split}
\end{equation}
In particular $b'=0$ wherever $b=0$.

\step{3: a first-order system where $b\ne0$}
On a component of $\{b\ne0\}$ choose a continuous complex square root
$y$ of $b$. Substituting $b=y^2$ in
\eqref{eq:obstruction-highest-mode} and dividing by $2y^4$,
\begin{equation}\label{eq:obstruction-square-root}
 (4y'-ay)\,y''+a\Bigl(y'-\frac yr\Bigr)^2
 +d\Bigl(yy'-\frac{y^2}r\Bigr)=0 .
\end{equation}
Fix a compact subinterval of $I$ and a constant $K_1>1+\sup|a|$, and
consider
\[
 U=\max\bigl(K_1|y|,\,|y'|\bigr)
   =\max\Bigl(K_1\sqrt{|b|},\ \frac{|b'|}{2\sqrt{|b|}}\Bigr).
\]
We claim $|U'|\le CU$ almost everywhere. Where $|y'|\ge K_1|y|$ we
have $|4y'-ay|\ge3|y'|$, and \eqref{eq:obstruction-square-root} gives
$|y''|\le C|y'|$. Where $K_1|y|$ is the larger term, its derivative is
at most $K_1|y'|\le K_1U$ in absolute value. The maximum of two
Lipschitz functions is Lipschitz, so the estimate holds across the
points where the maximum switches. Consequently $U$ is bounded above
and bounded away from zero on the component, up to its endpoints in
the compact subinterval.

\step{4: zeros are isolated and quadratic}
Let $r_0$ be an endpoint, interior to $I$, of a component of
$\{b\ne0\}$. There $y\to0$, so $U=|y'|$ nearby, and since
$|y''|\le C|y'|$ the derivative $y'$ has a limit $\lambda$, which is
not zero because $U$ is bounded below. Hence
\begin{equation}\label{eq:obstruction-quadratic-zero}
 b''=2\,(y')^2+2\,y\,y''\longrightarrow2\lambda^2\ne0
 \qquad(r\to r_0).
\end{equation}
So $r_0$ is an isolated zero of $b$ of order exactly two. Across such
a zero $U$ extends continuously, with the value $\sqrt{|b''|/2}$, and
satisfies the same differential inequality.

\step{5: continuation}
Consider the open set
$\mathcal O=\{b\ne0\}\cup\{b=0,\ b''\ne0\}$. Suppose a component of
$\mathcal O$ had a boundary point in the interior of $I$; there
$b=b''=0$. If zeros of $b$ accumulated at this point from inside the
component, the lower bound for $U$ at these quadratic zeros and the
continuity of $b''$ would give a contradiction; if a last interval
of $\{b\ne0\}$ ended at this point,
\eqref{eq:obstruction-quadratic-zero} would give a contradiction.
Hence $\mathcal O$ is empty or equal to $I$. By assumption $z^2$ is
radial on an open interval, so $b$ vanishes there and
$\mathcal O\ne I$. Therefore $b=0$ on $I$.

\step{6: conclusion}
Removing the highest coefficient and repeating the argument, all
coefficients $b_j$ with $j\ne0$ vanish, so $z^2$ is radial. For each
$r$ the function $\theta\mapsto z(r,\theta)$ is continuous on the
circle with constant square, hence constant.
\end{proof}

The proposition does not say that a single angular mode is useless: a
single mode is exactly the right seed for the linearized problem. It
says that the nonlinear correction of the seed cannot stay within any
finite number of modes once it has to meet a radial collar.
